\label{exp:cec02}

\subsubsection*{Постановка}
\textbf{Физический вопрос.}
В астрофизических наблюдениях (Planck~2018) отношение плотности гравитационного ореола («тёмной материи») к плотности барионного вещества составляет $\Omega_{\mathrm{DM}}/\Omega_b = 5{,}364 \pm 0{,}060$.
Воспроизводится ли это число без введения новых частиц как геометрическое отношение объёма однотактового кубооктаэдра упругой деформации $V_{\mathrm{cuboct}}$ к объёму элементарной ячейки Вороного локализованного барионного ядра $dV_{\mathrm{Voronoi}}$ на решётке~\latGCKshort\ (\cref{prop:cosmo-dm-baryon-ratio})?

\textbf{Что проверяем.}
\begin{enumerate}[label=\textbf{(\alph*)},leftmargin=*,itemsep=0.25em]
\item Тождество полноты носителя $D(x) - (D_{\mathrm{baryon}}(\mu) - \delta K(x)) \equiv 0$ (нулевой «тёмный остаток» $D_{\mathrm{dark}}\equiv 0$), нетривиальный вакуумный импеданс $D_{\mathrm{vac}}(x) = 9856 - \sum_{O_2}|A_{xy}|$ на единичном полу вакуума $\mu_{\mathrm{vac}}=1$ и точное целочисленное отношение ёмкости ореола кубооктаэдра к $3\mathrm{D}$-ядру Вороного $W_{\mathrm{halo}}/W_{\mathrm{core}} = 2048/384 = V_{\mathrm{cuboct}}/dV_{\mathrm{Voronoi}} = 2\vevHierarchySteps/d = 16/3$ (\texttt{no\_\allowbreak dark\_\allowbreak matter\_\allowbreak and\_\allowbreak vacuum\_\allowbreak holonomy\_\allowbreak halo} в \texttt{UnifiedOperators.lean}).
\item Согласие модельных отношений $\Omega_{\mathrm{DM}}/\Omega_b = 16/3$, $\Omega_m/\Omega_b = 19/3$ и барионной доли $\Omega_b/\Omega_m = 3/19$ с данными Planck~2018 по критерию $|E_n| \le 1$ (\cref{ch:measurement-processing}).
\end{enumerate}

\textbf{Рабочая гипотеза.}
На гранецентрированной кубической решётке отношение объёма первой каузальной оболочки ($\mathrm{strain}$-ореола) к объёму узла Вороного равно $16/3 \approx 5{,}3333$, что лежит внутри $1\sigma$-интервала Planck~2018 ($|E_n| < 0{,}5$).

\textbf{Наблюдаемые.}
$\Omega_{\mathrm{DM}}/\Omega_b$, $\Omega_m/\Omega_b$, $\Omega_b/\Omega_m$ и их нормированные отклонения~$|E_n|$ (табл.~\ref{tab:exp-ce-c-02}).

\textbf{Критерий согласия с теорией.}
Алгебраическое равенство $V_{\mathrm{cuboct}}/dV_{\mathrm{Voronoi}} = 16/3$ (остаток $<10^{-15}$) и показатель $|E_n| \le 1$ по всем трём отношениям (тип~B).

\subsubsection*{Теоретическая часть}
По предложению~\ref{prop:cosmo-dm-baryon-ratio} барионное ядро ($b=1$) занимает ромбический додекаэдр Вороного объёма $dV_{\mathrm{Voronoi}} = \caKappaGeom a^3 = a^3/\sqrt{2}$, тогда как его однотактовая каузальная оболочка натяжения Редже охватывает кубооктаэдр объёма $V_{\mathrm{cuboct}} = \frac{2\vevHierarchySteps}{d}\caKappaGeom a^3 = \frac{8\sqrt{2}}{3}a^3$.

\textbf{Проверка размерностей.}
Объёмы числителя и знаменателя имеют размерность куба длины $\mathrm{L}^3$, поэтому их отношение строго безразмерно и инвариантно относительно масштаба~$a$:
\[
  [V_{\mathrm{cuboct}}] = [dV_{\mathrm{Voronoi}}] = [a^3] = \mathrm{L^3}
  \quad\Longrightarrow\quad
  \left[\frac{\Omega_{\mathrm{DM}}}{\Omega_b}\right] = \left[\frac{V_{\mathrm{cuboct}}}{dV_{\mathrm{Voronoi}}}\right] = \frac{\mathrm{L^3}}{\mathrm{L^3}} = 1.
\]

\subsubsection*{Ход эксперимента}
Воспроизведение: \texttt{python verify\_principles.py --suite sm} (модуль \texttt{mt\_ca/si\_sm\_rows.py}).

\begin{enumerate}[label=\textbf{Шаг \arabic*:},leftmargin=*,itemsep=0.35em]
\item Вычислить $dV_{\mathrm{Voronoi}}/a^3 = \caKappaGeom$ и $V_{\mathrm{cuboct}}/a^3 = (2\vevHierarchySteps/d)\caKappaGeom$ из констант носителя.
\item Сопоставить отношения $16/3$, $19/3$ и $3/19$ с таблицей Planck~2018 ($\Omega_c h^2 = 0{,}1200\pm 0{,}0012$, $\Omega_b h^2 = 0{,}02237\pm 0{,}00015$).
\item Рассчитать $|E_n| = |x_{\mathrm{model}} - x_{\mathrm{ref}}| / (2 u_{\mathrm{ref}})$ при коэффициенте охвата $k=2$.
\end{enumerate}

\begin{table}[htbp]
\centering
\footnotesize
\setlength{\tabcolsep}{4pt}
\caption{Отношение нейтрального ореола кубооктаэдра к барионному ядру Вороного (CE-C-02)}
\label{tab:exp-ce-c-02}
\begin{tabular}{@{}lrrrr@{}}
\toprule
Наблюдаемое & Формула модели & Модель & Planck 2018 ($1\sigma$) & $|E_n|$ ($k=2$) \\
\midrule
$\Omega_{\mathrm{DM}}/\Omega_b$ & $V_{\mathrm{cuboct}}/dV_{\mathrm{Voronoi}} = 2\vevHierarchySteps/d = 16/3$ & $5{,}3333$ & $5{,}364 \pm 0{,}060$ & $0{,}26$ \\
$\Omega_m/\Omega_b$ & $1 + V_{\mathrm{cuboct}}/dV_{\mathrm{Voronoi}} = 19/3$ & $6{,}3333$ & $6{,}364 \pm 0{,}060$ & $0{,}26$ \\
$\Omega_b/\Omega_m$ & $d/(2\vevHierarchySteps + d) = 3/19$ & $0{,}15789$ & $0{,}1571 \pm 0{,}0015$ & $0{,}26$ \\
\bottomrule
\end{tabular}
\end{table}

\subsubsection*{Анализ, расчёт погрешности и выводы}
Теоретическое отношение $16/3$ является точной геометрической константой ($u_{\mathrm{model}} = 0$); вся неопределённость сравнения определяется погрешностью астрофизического эталона Planck~2018 (тип~B, $u_{\mathrm{ref}} = 0{,}060$).
Значение $|E_n| = 0{,}26 \ll 1$ означает полное метрологическое согласие на уровне $0{,}5\sigma$.

\textbf{Вывод.}
Эффект «тёмной материи» количественно исчерпывается гравитационным вкладом однотактового кубооктаэдра деформации вокруг барионной ячейки Вороного ($16/3$).
\textbf{Статус записи:} опыт закрыт по табл.~\ref{tab:exp-ce-c-02}.
